
%%%%%%%%%%
%% 2009 %%
%%%%%%%%%%


@article{HendersonMalani2009corporatePhilanthropy,
   author = {Henderson, M. Todd and Malani, Anup},
   title = {Corporate Philanthropy and the Market for Altruism Essay},
   journal = {Columbia Law Review},
   volume = {109},
   number = {3},
   pages = {571-628},
   url = {https://heinonline.org/HOL/P?h=hein.journals/clr109&i=583},
   year = {2009},
   type = {Journal Article}
}

@article{LaxminarayanMalani2009abxResistanceIncentives,
   author = {Laxminarayan, R. and Malani, A.},
   title = {The right combination of carrots and sticks},
   journal = {Resources},
   number = {173},
   pages = {13–15},
   abstract = {Individual countries reporting an outbreak would face trade and economic sanctions, and therefore, had every incentive to remain quiet. This article discusses incentives for countries to look for and report disease outbreaks as a means to control epidemics.},
   ISSN = {0048-7376},
   year = {2009},
   type = {Journal Article}
}

%%%%%%%%%%
%% 2008 %%
%%%%%%%%%%


@article{Malani2008nonProfitSignalQuality,
   author = {Malani, Anup and David, Guy},
   title = {Does Nonprofit Status Signal Quality?},
   journal = {The Journal of Legal Studies},
   volume = {37},
   number = {2},
   pages = {551-576},
   note = {doi: 10.1086/589668},
   abstract = {Abstract A popular theory for why firms take nonprofit status is that it is a signal of quality. This paper offers a simple, empirical test of this theory. If nonprofit status signals quality, surely nonprofit firms would want to ensure that consumers are aware of this. A natural way for firms to do this is to indicate their nonprofit status in their advertising. Taking this cue, we conducted a survey of over 2,800 firms in the hospital, nursing home, or child care industries in order to determine whether nonprofit firms communicate their status to consumers on their Web sites or yellow pages listings. We find that fewer than 7.5 percent of nonprofit firms signal their status in yellow pages listings, only 25 percent do so on their home pages, and 30 percent do so on their about?us pages. Indeed, over 35 percent never signal their nonprofit status on their Web sites. Our evidence does not support the hypothesis that nonprofit status is a signal of quality.
A popular theory for why firms take nonprofit status is that it is a signal of quality. This paper offers a simple, empirical test of this theory. If nonprofit status signals quality, surely nonprofit firms would want to ensure that consumers are aware of this. A natural way for firms to do this is to indicate their nonprofit status in their advertising. Taking this cue, we conducted a survey of over 2,800 firms in the hospital, nursing home, or child care industries in order to determine whether nonprofit firms communicate their status to consumers on their Web sites or yellow pages listings. We find that fewer than 7.5 percent of nonprofit firms signal their status in yellow pages listings, only 25 percent do so on their home pages, and 30 percent do so on their about?us pages. Indeed, over 35 percent never signal their nonprofit status on their Web sites. Our evidence does not support the hypothesis that nonprofit status is a signal of quality.},
   ISSN = {0047-2530},
   DOI = {10.1086/589668},
   url = {https://doi.org/10.1086/589668},
   year = {2008},
   type = {Journal Article}
}

@article{Malani2008regulationPlaceboEffects,
   author = {Malani, Anup},
   title = {Regulation with Placebo Effects},
   journal = {Duke Law Journal},
   volume = {58},
   number = {3},
   pages = {411-472},
   url = {https://heinonline.org/HOL/P?h=hein.journals/duklr58&i=417},
   year = {2008},
   type = {Journal Article}
}

@article{Malani2008caffeinePlaceboEffects,
   author = {Malani, Anup and Houser, Daniel},
   title = {Expectations mediate objective physiological placebo effects},
   journal = {Advances in health economics and health services research},
   volume = {20},
   pages = {311-327},
   abstract = {&lt;h4>Purpose&lt;/h4>A placebo effect is a (positive) change in health outcomes that is due to a (positive) change in beliefs about the value of a treatment. Placebo effects might be "behavioral," in the sense that revised beliefs lead to behavioral changes or new actions that in turn yield changes in health outcomes. Placebo effects might also include a "physiological" component, which refers broadly to non-behavioral, brain-modulated mechanisms by which new beliefs cause changes in health outcomes. Nearly all formal economic models of human behavior are consistent with behavioral placebo effects, but strongly inconsistent with their physiological counterparts. The reason is that the latter effects can imply that expectations enter, rather than multiply, state-contingent preferences. It is therefore unfortunate that little evidence exists on physiological placebo effects. We report data from novel clinical experiments with caffeine that seek to provide such evidence.&lt;h4>Methods&lt;/h4>Subjects visit the clinic on multiple occasions. On each visit they ingest either a placebo or caffeine pill. Subjects only know the probability with which the pill includes caffeine. We obtain physiological measurements prior to ingestion and at 30, 60, and 90 min after ingestion. Importantly, we constrain subjects to remain seated and read preselected magazines during the interval between treatment and outcome measurement.&lt;h4>Findings&lt;/h4>Our design provides particularly clean inference because it (i) eliminates the possibility of behavioral confounds; (ii) provides for measurements at the individual level; (iii) manipulates beliefs without deception; and (iv) uses salient rewards. We find evidence for the existence of physiological placebo effects mediated by expectations.&lt;h4>Implications&lt;/h4>Our results are consistent with the possibility that the prefrontal cortex provides external, top-down control that modulates physiological outcomes, and make a case for the importance of research geared toward developing appropriate and tractable frameworks that accommodate non-linear relationships between expectations and preferences.},
   keywords = {Health Knowledge, Attitudes, Practice
Placebo Effect},
   ISSN = {0731-2199},
   url = {http://europepmc.org/abstract/MED/19552314},
   year = {2008},
   type = {Journal Article}
}

@article{Malani2008patentEnrollmentTrials,
   author = {Malani, Anup},
   title = {Patient enrollment in medical trials: Selection bias in a randomized experiment},
   journal = {Journal of Econometrics},
   volume = {144},
   number = {2},
   pages = {341-351},
   abstract = {Self-selection can bias estimates of treatment effects from randomized experiments if one is interested in extrapolating results to the general population. This paper notes that there is an isomorphism between the Roy model for the sorting of workers into sectors and the decision of subjects to participate in randomized experiments. The main implication is that, as the probability of receiving active treatment rises, patients who are less optimistic about new treatment will begin to enroll and estimates of treatment effects will fall. This, in turn, implies that selection bias is positive. These findings are confirmed with data from trials of ulcer medications.},
   keywords = {Selection
Randomization bias
Experiment
Treatment effect},
   ISSN = {0304-4076},
   DOI = {https://doi.org/10.1016/j.jeconom.2008.04.002},
   url = {https://www.sciencedirect.com/science/article/pii/S0304407608000328},
   year = {2008},
   type = {Journal Article}
}

@article{BuchmanMalani2008socialNormsAbxOveruse,
   author = {Buchman, Timothy G  and Dushoff, Jonathan  and Effron, Mark B  and Ehrlich, Paul R  and Fitzpatrick, Susan  and Laxminarayan, Ramanan  and Levin, Bruce  and Levin, Simon A  and Lipsitch, Marc  and Malani, Anup  and Nemeroff, Carol  and Otto, Sarah P  and Patel, Vimla L  and Solomkin, Joseph S },
   title = {Antibiotic Overuse: The Influence of Social Norms},
   journal = {Journal of the American College of Surgeons},
   volume = {207},
   number = {2},
   ISSN = {1879-1190},
   url = {https://journals.lww.com/journalacs/fulltext/2008/08000/antibiotic_overuse__the_influence_of_social_norms.17.aspx},
   year = {2008},
   type = {Journal Article}
}

@article{Malani2008lawAsLocalAmenity,
   author = {Malani, Anup},
   title = {Valuing Laws as Local Amenities},
   journal = {Harvard Law Review},
   volume = {121},
   number = {5},
   pages = {1273-1331},
   url = {https://heinonline.org/HOL/P?h=hein.journals/hlr121&i=1295},
   year = {2008},
   type = {Journal Article}
}

% NON-JOURNAL

@inproceedings{Malani2008conferenceCommentHealthInsurance,
   author = {Malani, Anup},
   title = {Comment on Chapter 3 (Burman, Garrett and Khitatrakun)},
   booktitle = {Conference on Taxes and Health Insurance: Analysis and Policy},
   editor = {Burman, Leonard E. and Helms, Robert B.},
   publisher = {Brookings Institution Press},
   year = {2008},
   type = {Conference Proceedings}
}

%%%%%%%%%%
%% 2007 %%
%%%%%%%%%%

@article{MalaniPosner2007forProfitCharities,
   author = {Malani, Anup and Posner, Eric A.},
   title = {The Case for For-Profit Charities},
   journal = {Virginia Law Review},
   volume = {93},
   number = {8},
   pages = {2017-2067},
   ISSN = {00426601},
   url = {http://www.jstor.org/stable/25050404},
   year = {2007},
   type = {Journal Article}
}



@article{Malani2007medMalReview,
   author = {Malani, Anup},
   title = {Introduction: Current Research on Medical Malpractice Liability},
   journal = {The Journal of Legal Studies},
   volume = {36},
   number = {S2},
   pages = {S1-S8},
   note = {doi: 10.1086/587430},
   ISSN = {0047-2530},
   DOI = {10.1086/587430},
   url = {https://doi.org/10.1086/587430},
   year = {2007},
   type = {Journal Article}
}

% NON-JOURNAL

@book{LaxminarayanMalani2007extendingTheCureAbxResistance,
   author = {Laxminarayan, R. and Malani, A.},
   title = {Extending the Cure: Policy Responses to the Growing Threat of Antibiotic Resistance},
   publisher = {Resources for the Future},
   ISBN = {9781933115573},
   url = {https://books.google.com/books?id=xTgD4L-u5OEC},
   year = {2007},
   type = {Book}
}

% WORKING PAPER

@inproceedings{RN6614,
   author = {Beylin, Ilya and Malani, Anup and Abrams, David},
   title = {Finding Love in the Wreckage: Estimating Spousal Altruism with Data on Fatal Car Accidents},
   booktitle = {3rd Annual Conference on Empirical Legal Studies Papers},
   url = {https://papers.ssrn.com/sol3/papers.cfm?abstract_id=1119379},
   type = {Conference Proceedings}
}

%%%%%%%%%%
%% 2006 %%
%%%%%%%%%%

@article{Malani2006habeasSettlements,
@article{RN6579,
   author = {Malani, Anup},
   title = {Habeas Settlements},
   journal = {Virginia Law Review},
   volume = {92},
   number = {1},
   pages = {1-68},
   url = {https://heinonline.org/HOL/P?h=hein.journals/valr92&i=15},
   year = {2006},
   type = {Journal Article}
}

@article{Malani2006JPEplaceboEffects,
   author = {Malani, Anup},
   title = {Identifying Placebo Effects with Data from Clinical Trials},
   journal = {Journal of Political Economy},
   volume = {114},
   number = {2},
   pages = {236-256},
   note = {doi: 10.1086/500279},
   abstract = {A medical treatment is said to have placebo effects if patients who are optimistic about the treatment respond better to the treatment. This paper proposes a simple test for placebo effects. Instead of comparing the treatment and control arms of a single trial, one should compare the treatment arms of two trials with different probabilities of assignment to treatment. If there are placebo effects, patients in the higher?probability trial will experience better outcomes simply because they believe that there is a greater chance of receiving treatment. This paper finds evidence of placebo effects in trials of antiulcer and cholesterol?lowering drugs.},
   ISSN = {0022-3808},
   DOI = {10.1086/500279},
   url = {https://doi.org/10.1086/500279},
   year = {2006},
   type = {Journal Article}
}

%%%%%%%%%%
%% 2005 %%
%%%%%%%%%%


% WORKING PAPERS

@article{MalaniMullin2005reallocationJointSeveral,
   author = {Malani, Anup and Mullin, Charles},
   title = {Assessing the Merits of Reallocation under Joint and Several Liability: Evidence from Asbestos Litigation},
   journal = {University of Virginia Law School Working Paper Series},
   url = {https://www.researchgate.net/profile/Charles-Mullin-2/publication/228435323_Assessing_the_Merits_of_Reallocation_under_Joint_and_Several_Liability_Evidence_from_Asbestos_Litigation/links/0c960537e04aec09a6000000/Assessing-the-Merits-of-Reallocation-under-Joint-and-Several-Liability-Evidence-from-Asbestos-Litigation.pdf},
   year = {2005},
   type = {Journal Article}
}



%%%%%%%%%%
%% 2004 %%
%%%%%%%%%%

@article{RN6618,
   author = {Hynes, Richard M and Malani, Anup and Posner, Eric A},
   title = {The Political Economy of Property Exemption Laws},
   journal = {The Journal of Law and Economics},
   volume = {47},
   number = {1},
   pages = {19-43},
   note = {doi: 10.1086/386276},
   abstract = {Abstract Exemption laws enable people who default on loans to protect certain assets from liquidation. Every state has its own set of exemption laws, and they vary widely. The 1978 federal bankruptcy law contains a set of national exemptions, which debtors in bankruptcy are permitted to use instead of their state?s exemptions unless the state has formally ?opted out? of the federal system. We contend that states? decisions to opt out shed light on their exemption levels. We find that states are more likely to opt out if their state exemption is lower than the federal exemption and that states are more likely to opt out if they also have a high bankruptcy filing rate and transfer little money to the poor. These latter findings suggest that studies that examine the impact of exemptions on, for example, the bankruptcy rate should not treat exemption levels as exogenous variables.
Exemption laws enable people who default on loans to protect certain assets from liquidation. Every state has its own set of exemption laws, and they vary widely. The 1978 federal bankruptcy law contains a set of national exemptions, which debtors in bankruptcy are permitted to use instead of their state?s exemptions unless the state has formally ?opted out? of the federal system. We contend that states? decisions to opt out shed light on their exemption levels. We find that states are more likely to opt out if their state exemption is lower than the federal exemption and that states are more likely to opt out if they also have a high bankruptcy filing rate and transfer little money to the poor. These latter findings suggest that studies that examine the impact of exemptions on, for example, the bankruptcy rate should not treat exemption levels as exogenous variables.},
   ISSN = {0022-2186},
   DOI = {10.1086/386276},
   url = {https://doi.org/10.1086/386276},
   year = {2004},
   type = {Journal Article}
}

% LETTER

@article{CantorMalani2004placeboEffects,
   author = {Cantor, Harvey E. and Malani, Anup},
   title = {Ethical Issues in Research in Complementary and Alternative Medicine},
   journal = {JAMA},
   volume = {291},
   number = {18},
   pages = {2192-2194},
   ISSN = {0098-7484},
   DOI = {10.1001/jama.291.18.2192-a},
   url = {https://doi.org/10.1001/jama.291.18.2192-a},
   year = {2004},
   type = {Journal Article}
}

% WORKING PAPERS

@article{Malani2004optionValueTherapeutics,
   author = {Malani, Anup and Hu, Feifang},
   title = {The option value of new therapeutics},
   journal = {Available at SSRN 617382},
   url = {https://papers.ssrn.com/sol3/papers.cfm?abstract_id=617382},
   year = {2004},
   type = {Journal Article}
}

@article{MalaniChoi2004nonProfitNursingHomes,
   author = {Malani, Anup and Choi, Albert H},
   title = {Are non-profit firms simply for-profits in disguise? Evidence from executive compensation in the nursing home industry},
   journal = {Evidence from Executive Compensation in the Nursing Home Industry (September 26, 2004)},
   url = {https://papers.ssrn.com/sol3/papers.cfm?abstract_id=617362},
   year = {2004},
   type = {Journal Article}
}

%%%%%%%%%%
%% 2003 %%
%%%%%%%%%%

% BOOK

@inbook{Malani2003nonProfit,
   author = {Malani, Anup and Philipson, Tomas and David, Guy},
   title = {Theories of firm behavior in the nonprofit sector. A synthesis and empirical evaluation},
   booktitle = {The governance of not-for-profit organizations},
   editor = {Glaeser, Edward},
   publisher = {NBER/University of Chicago Press},
   pages = {181-216},
   year = {2003},
   type = {Book Section}
}

% THESIS

@phdthesis{Malani2003placeboEffects,
   author = {Malani, Anup},
   title = {Placebo effects, self-selection and the external validity of clinical trials},
   url = {https://www.proquest.com/pqdtglobal/docview/305291636/B27E4F5857104EBAPQ/4?accountid=14657&sourcetype=Dissertations%20&%20Theses},
   year = {2003},
   type = {Thesis}
}

%%%%%%%%%%
%% 2002 %%
%%%%%%%%%%

% WORKING PAPERS

@article{Malani2002felonyMurder,
   author = {Malani, Anup},
   title = {Does the felony murder rule deter? Evidence from FBI crime data},
   journal = {Working Paper},
   url = {https://www.researchgate.net/profile/Anup-Malani/publication/265435387_Does_the_Felony-Murder_Rule_Deter_Evidence_from_FBI_Crime_Data/links/5745b88708ae298602f9d20f/Does-the-Felony-Murder-Rule-Deter-Evidence-from-FBI-Crime-Data.pdf},
   year = {2002},
   type = {Journal Article}
}


